% GENERATED FILE. Edit canonical lessons, then run npm run generate:books.
% canonical-source-hash: fnv1a64:f6dd0b571a736324
% canonical-lessons: FA-C32-narahat, FA-C32-pezeshk, FA-C32-pedarbozorg, FA-C32-madarbozorg, FA-C32-dayi

\chapter{Upset, Doctor, Grandfather, Grandmother, Uncle}
\label{ch:fa-upset-doctor-grandfather-grandmother-uncle}
\hlchaptermodality{32}

\begin{chapteropening}
\textbf{By the end of this chapter:} \emph{I can say upset, sad, a doctor, a grandfather, a grandmother and an uncle.}

Complete the last lesson of chapter 32: I can say upset, sad, a doctor, a grandfather, a grandmother and an uncle.
\end{chapteropening}

\section[\texorpdfstring{\fa{ناراحت} (nârâhat)}{ناراحت (nârâhat)}]{\fa{ناراحت} (nârâhat) — upset, sad}

\label{lesson:FA-C32-narahat}



\begin{quote}
Before the new one: say the Persian for hungry, then the Persian for happy.
\end{quote}

\subsection*{You'll want to know: \fa{ناراحت}}
\textbf{\fa{ناراحت}} — \emph{nârâhat} — \textquotedblleft{}upset, sad\textquotedblright{}.

\begin{cousinweb}
\emph{nârâhat} is \emph{nâ-}, \textquotedblleft{}not\textquotedblright{}, with \emph{râhat}, \textquotedblleft{}ease, comfort\textquotedblright{}: \textquotedblleft{}not at ease\textquotedblright{}. \emph{nâ-} is the native Persian negative prefix; \emph{râhat} is borrowed from Arabic \emph{rāḥa}, \textquotedblleft{}rest, ease\textquotedblright{} (root \emph{r-w-ḥ}, Lane).
\end{cousinweb}

\begin{grammarlens}[title={What you've built}]
One more word for this chapter.
\end{grammarlens}

\subsection*{Your turn}
\begin{itemize}
\raggedright
  \item \emph{Say it:} \emph{nârâhat}
  \item \emph{Say it:} \emph{nârâhat}, once more
  \item \emph{Say it:} say \emph{nârâhat}
  \item \emph{Recall:} say the Persian for hungry, then the Persian for happy, then say \emph{nârâhat} again
\end{itemize}

\begin{tcolorbox}[breakable,colback=teal!4,colframe=teal!35!black,title={Before you move on}]
What does \fa{ناراحت} mean? (\textquotedblleft{}Upset, sad\textquotedblright{}.) Say it once more.
\end{tcolorbox}
\section[\texorpdfstring{\fa{پزشک} (pezeshk)}{پزشک (pezeshk)}]{\fa{پزشک} (pezeshk) — a doctor}

\label{lesson:FA-C32-pezeshk}



\begin{quote}
Before the new one: say the Persian for happy, then the Persian for upset.
\end{quote}

\subsection*{You'll want to know: \fa{پزشک}}
\textbf{\fa{پزشک}} — \emph{pezeshk} — \textquotedblleft{}a doctor\textquotedblright{}.

\begin{cousinweb}
\emph{pezeshk} is a native Persian word, inherited rather than borrowed; it is usually traced to Middle Persian \emph{bizešk}, \textquotedblleft{}physician\textquotedblright{}.
\end{cousinweb}

\begin{grammarlens}[title={What you've built}]
One more word for this chapter.
\end{grammarlens}

\subsection*{Your turn}
\begin{itemize}
\raggedright
  \item \emph{Say it:} \emph{pezeshk}
  \item \emph{Say it:} \emph{pezeshk}, once more
  \item \emph{Say it:} say \emph{pezeshk}
  \item \emph{Recall:} say the Persian for happy, then the Persian for upset, then say \emph{pezeshk} again
\end{itemize}

\begin{tcolorbox}[breakable,colback=teal!4,colframe=teal!35!black,title={Before you move on}]
What does \fa{پزشک} mean? (\textquotedblleft{}A doctor\textquotedblright{}.) Say it once more.
\end{tcolorbox}
\section[\texorpdfstring{\fa{پدربزرگ} (pedarbozorg)}{پدربزرگ (pedarbozorg)}]{\fa{پدربزرگ} (pedarbozorg) — a grandfather}

\label{lesson:FA-C32-pedarbozorg}



\begin{quote}
Before the new one: say the Persian for upset, then the Persian for a doctor.
\end{quote}

\subsection*{You'll want to know: \fa{پدربزرگ}}
\textbf{\fa{پدربزرگ}} — \emph{pedarbozorg} — \textquotedblleft{}a grandfather\textquotedblright{}.

\begin{cousinweb}
\emph{pedarbozorg} is built from two inherited Persian words: \emph{pedar}, \textquotedblleft{}father\textquotedblright{} (Middle Persian \emph{pidar}), and \emph{bozorg}, \textquotedblleft{}big, great\textquotedblright{} (Middle Persian \emph{wuzurg}, EtymDB). The \textquotedblleft{}great father\textquotedblright{} is the grandfather.
\end{cousinweb}

\begin{grammarlens}[title={What you've built}]
One more word for this chapter.
\end{grammarlens}

\subsection*{Your turn}
\begin{itemize}
\raggedright
  \item \emph{Say it:} \emph{pedarbozorg}
  \item \emph{Say it:} \emph{pedarbozorg}, once more
  \item \emph{Say it:} say \emph{pedarbozorg}
  \item \emph{Recall:} say the Persian for upset, then the Persian for a doctor, then say \emph{pedarbozorg} again
\end{itemize}

\begin{tcolorbox}[breakable,colback=teal!4,colframe=teal!35!black,title={Before you move on}]
What does \fa{پدربزرگ} mean? (\textquotedblleft{}A grandfather\textquotedblright{}.) Say it once more.
\end{tcolorbox}
\section[\texorpdfstring{\fa{مادربزرگ} (mâdarbozorg)}{مادربزرگ (mâdarbozorg)}]{\fa{مادربزرگ} (mâdarbozorg) — a grandmother}

\label{lesson:FA-C32-madarbozorg}



\begin{quote}
Before the new one: say the Persian for a doctor, then the Persian for a grandfather.
\end{quote}

\subsection*{You'll want to know: \fa{مادربزرگ}}
\textbf{\fa{مادربزرگ}} — \emph{mâdarbozorg} — \textquotedblleft{}a grandmother\textquotedblright{}.

\begin{cousinweb}
\emph{mâdarbozorg} is built from two inherited Persian words: \emph{mâdar}, \textquotedblleft{}mother\textquotedblright{} (Middle Persian \emph{mādar}), and \emph{bozorg}, \textquotedblleft{}big, great\textquotedblright{} (Middle Persian \emph{wuzurg}, EtymDB). The \textquotedblleft{}great mother\textquotedblright{} is the grandmother.
\end{cousinweb}

\begin{grammarlens}[title={What you've built}]
One more word for this chapter.
\end{grammarlens}

\subsection*{Your turn}
\begin{itemize}
\raggedright
  \item \emph{Say it:} \emph{mâdarbozorg}
  \item \emph{Say it:} \emph{mâdarbozorg}, once more
  \item \emph{Say it:} say \emph{mâdarbozorg}
  \item \emph{Recall:} say the Persian for a doctor, then the Persian for a grandfather, then say \emph{mâdarbozorg} again
\end{itemize}

\begin{tcolorbox}[breakable,colback=teal!4,colframe=teal!35!black,title={Before you move on}]
What does \fa{مادربزرگ} mean? (\textquotedblleft{}A grandmother\textquotedblright{}.) Say it once more.
\end{tcolorbox}
\section[\texorpdfstring{\fa{دایی} (dâyi)}{دایی (dâyi)}]{\fa{دایی} (dâyi) — an uncle (mother's brother)}

\label{lesson:FA-C32-dayi}



\begin{quote}
Before the new one: say the Persian for a grandfather, then the Persian for a grandmother.
\end{quote}

\subsection*{You'll want to know: \fa{دایی}}
\textbf{\fa{دایی}} — \emph{dâyi} — \textquotedblleft{}an uncle (mother's brother)\textquotedblright{}.

\begin{cousinweb}
\emph{dâyi} is usually traced to Turkic, to the word behind Turkish \emph{dayı}, \textquotedblleft{}maternal uncle\textquotedblright{}; EtymDB instead connects it with Persian \emph{dâ}, \textquotedblleft{}mother\textquotedblright{}. The origin is not settled.
\end{cousinweb}

\begin{grammarlens}[title={What you've built}]
One more word for this chapter.
\end{grammarlens}

\subsection*{Your turn}
\begin{itemize}
\raggedright
  \item \emph{Say it:} \emph{dâyi}
  \item \emph{Say it:} \emph{dâyi}, once more
  \item \emph{Say it:} say \emph{dâyi}
  \item \emph{Recall:} say the Persian for a grandfather, then the Persian for a grandmother, then say \emph{dâyi} again
\end{itemize}

\begin{tcolorbox}[breakable,colback=teal!4,colframe=teal!35!black,title={Before you move on}]
What does \fa{دایی} mean? (\textquotedblleft{}An uncle (mother's brother)\textquotedblright{}.) Say it once more.
\end{tcolorbox}
